\documentclass[12pt]{article}
\usepackage[utf8]{inputenc}
\usepackage[T1]{fontenc}
\usepackage{hyperref}
\title{Yaminahua Language Resources for K-12 Curriculum Development}
\author{Compiled by Claude Sonnet 5.5 (claude-sonnet-5-5) for Andr\'as Kornai}
\date{}
\begin{document}
\maketitle

\section{Introduction}
Yaminahua (ISO 639-3 yaa; Panoan) is spoken in Peru, Bolivia and Brazil. Speaker counts in the results vary widely, from under 1,400 for Peru in one source to about 3,000 in others, probably because they cover different areas. The search found one grammar, a documentation archive, an official alphabet and mentions of bilingual schools. It found no dictionary and no textbooks.

\section{Primary Educational Resources}
\begin{itemize}
\item A 2017 ministerial resolution is reported to have set an official 21-letter alphabet for Yaminahua. Intercultural bilingual schools are reported (five registered as of 2013, including students from other peoples): \url{https://www.ecured.cu/Lengua_yaminawa}
\item Omniglot says the language is taught in some schools but that few speakers read and write it: \url{https://omniglot.com/writing/yaminawa.htm}
\end{itemize}

\section{Linguistic and Reference Materials}
\begin{itemize}
\item Faust and Loos (2002), \emph{Gram\'atica del idioma yaminahua} (SIL, Lima): \url{https://ldh.clld.org/sources/item_402694}
\item SAPhon phonological inventory: \url{https://linguistics.berkeley.edu/saphon/en/inv/Yaminahua.html}
\end{itemize}

\section{Religious and Community Resources}
None were found.

\section{Digital and Audio Resources}
\begin{itemize}
\item ``Multidialectal Lexical Documentation of Yaminahua, Nahua, and Sharanahua'' (ELDP project, archived at ELAR): \url{https://cultureincrisis.org/projects/multidialectal-lexical-documentation-of-yaminahua-nahua-and-sharanahua}
\item Archive record: \url{https://eldp.access.preservica.com/dk0605/}
\end{itemize}

\section{Curriculum Development Context}
Schooling in the language appears limited. The official alphabet and the SIL grammar give a starting point for orthography and grammar. Primers, a dictionary and Ministry of Education materials were not found and should be sought, and the ELAR lexical documentation could support a dictionary. Development should be done with Yaminahua communities and teachers.
\end{document}